\documentclass[11pt,letterpaper]{article}
\usepackage[T1]{fontenc}
\usepackage{lmodern}
\usepackage[margin=1in]{geometry}
\usepackage{amsmath,amssymb,amsthm,mathtools,microtype}
\usepackage{booktabs,array,enumitem,xurl,xcolor}
\usepackage[colorlinks=true,linkcolor=blue!50!black,citecolor=blue!50!black,urlcolor=blue!50!black]{hyperref}
\hypersetup{pdftitle={Near-extremal entropy for admissible Collatz words},pdfauthor={Jake S. Brusca}}
\pdfinfoomitdate=1
\pdftrailerid{}
\pdfsuppressptexinfo=15
\setlength{\emergencystretch}{1.5em}
\setlength{\parskip}{0.25em}
\newtheorem{theorem}{Theorem}
\title{Near-extremal entropy for admissible Collatz words\\[0.3em]\large A written rung-limit and large-tilt research argument}
\author{Jake S. Brusca}
\date{28 November 2026}
\begin{document}
\maketitle
\begin{abstract}
We develop a written near-maximal reward argument for admissible exponent
words in the model of Polson and Sokolov. An exact endpoint transformation,
a bounded run-cutting cost and independent finite-block counting give a
candidate near-one rung entropy limit. A large-tilt sandwich yields the
proposed leading and second-order R9 expansion. The argument supplies only a
near-one interpretation of M4 and convergence for large tilt in R8, not the
full problems. No end-to-end formal proof, independent kernel replay or human
expert refereeing of these infinite rung claims is certified here. A separate
appendix proves nonconvergence for the fixed-reference all-length extension
at c=1; it does not replace the intended rung-only M4 question.
\end{abstract}
\begin{quote}\small
\textbf{Personal project; work in progress.} I am not an expert in this field.
I am exploring these ideas for fun and to learn. Corrections, counterexamples
and suggestions are welcome. ChatGPT assisted with ideas, proofs, code,
writing and review. AI review is not human expert refereeing, acceptance or a
priority claim. No ordinary Collatz convergence or global forced capture follows.
\end{quote}
\textbf{Evidence status.} The argument below is a written proof candidate, reviewed in a separate AI context. It has not been refereed by a human expert, formalized end to end, or independently kernel checked. The public finite bridge library proves the literal finite dictionary, unbounded rungs and unique maximum. The separate all-length appendix has its own 23 frozen formal statements and pending independent assurance. Neither package proves the main rung-pressure, entropy, critical regularity or R9 claims below. Finite exact tests check model translations; they do not establish limits. No claim of priority or external acceptance is made.

The intended contribution is a rung entropy limit near maximal reward and its near-extremal expansion. It addresses the rung interpretation of M4 near one, the large-tilt convergence part of R8, and the proposed asymptotic in R9 of Polson and Sokolov. It does not prove the complete R8 question, compute the requested M4 entropy table, establish an unrestricted M4 limit, or imply anything about integer Collatz cycles or convergence.

Primary source: Nicholas G. Polson and Vadim Sokolov~\cite{polson}, \emph{On the Collatz Conjecture: Martingales, Momentum and Continued Fractions}, edition 1 (2026), \href{https://github.com/VadimSokolov/collatz/blob/5840497a23fc0d64127b850aea68862d7c93f2b7/book/chapters/ch10h_ratchets.tex}{v1.4.0 ratchets chapter}, labels \texttt{rrt:sec:m4}, \texttt{rrt:prob:r8}, \texttt{rrt:prob:r9}; the tag resolves to commit \nolinkurl{5840497a23fc0d64127b850aea68862d7c93f2b7}. The \href{https://github.com/VadimSokolov/collatz/blob/5840497a23fc0d64127b850aea68862d7c93f2b7/book/chapters/ch10f_missed.tex}{missed-problems chapter} uses unrestricted-limit wording for M4; the detailed ratchets discussion fixes the rung subsequence. These formulations must be distinguished. The \href{../external_problems/README.md}{status register} distinguishes the exact source populations and records the outstanding assurance gates.

\section{Related work and the use of classical methods}
Mechanical extremizers for rational Collatz cycles predate this draft.
Halbeisen and Hungerb\"uhler~\cite{halbeisen-hungerbuhler}, Lemma~5 and
Corollary~1, identify balanced binary words in a maximum-minimum cycle problem;
their equation~(4) gives the affine concatenation identity. Their binary
length/odd-count convention must be converted to our stride-length/halving-count
convention. That extremal result does not count the near-maximal words studied
here; no new priority for mechanical extremizers is claimed.

The phase averages use classical irrational-rotation equidistribution
\cite{weyl}. The pressure argument uses the familiar constrained-partition
concatenation mechanism of renewal and pinning models; compare
Giacomin~\cite{giacomin}, Sections~1.2 and~3.1, especially equation~(3.1).
Here the environment is a deterministic irrational rotation. Independent
random disorder, almost-sure pressure limits and critical regularity from
that literature do not automatically transfer to this model.

The concavity and duality language follows standard convex analysis
\cite{rockafellar}. The matching microcanonical lower bound below still needs
its explicit finite-block construction. The endpoint Laplace calculation is
an instance of Watson's lemma~\cite{dlmf}, Section~2.3(ii); the polymer error
bounds and resulting constants remain specific claims of this draft. These
comparisons identify prior methods, without certifying the proposed entropy
limit, critical differentiability or novelty.

\section{Definitions and the proposed statement}

All logarithms below are natural unless a base is displayed. Set
\[
\theta=\log_2 3,\qquad \beta=\theta-1,\qquad
\delta(r)=\lceil r\theta\rceil-r\theta.
\]

A \textbf{rung} is a positive integer r such that $\delta(r)<\delta(t)$ for every positive integer t<r. At length r, a source word is a tuple of positive integers $k_0,\ldots,k_{r-1}$ with prefix sums $K_t=\sum_{i<t}k_i$, endpoint $K_r=\lceil r\theta\rceil$, and caps
\[
K_t\le \lfloor t\theta+\delta(r)\rfloor\quad(0\le t<r).
\]

Its reward and the phase-zero reference are
\[
\Phi(k)=\sum_{t=0}^{r-1}2^{K_t-t\theta},\qquad
B_r=\sum_{t=0}^{r-1}2^{-\{t\theta\}}.
\]

Let N(c,r) count these words with $\Phi(k)\ge cB_r$. The candidate result is: for $0<\epsilon<2(\log 2)10^{-4}$, the limit
\[
\rho(1-\epsilon)=\lim_{\substack{r\to\infty\\r\text{ a rung}}}
 \frac{\log_{10}N(1-\epsilon,r)}r
\]

exists, and, as $\epsilon\downarrow0$ \textbf{after this length limit},
\[
\rho(1-\epsilon)=\frac{3\epsilon}{2\log 2\log 10}
 \left[\log(1/\epsilon)-\log\log(1/\epsilon)
       +1+\log(2/3)+o(1)\right].
\]

At $c=1$ the rung population has exactly one word, so $\rho(1)=0$. The displayed leading coefficient is approximately 0.93983174299; this decimal is illustrative, not an independently certified numerical enclosure.

We use the following classical facts, with the stated hypotheses: an irrational rotation has uniformly convergent averages, in its initial phase, for bounded Riemann-integrable functions; finite independent bounded random variables obey Chebyshev's inequality; a finite concave function is continuous on the interior of its interval; and elementary calculus and the geometric series. There is no assumption of independent symbols along an integer Collatz orbit.

\section{Exact rung and depth transformations}

The number $\theta$ is irrational: a rational identity $\theta=a/b$ would give $2^a=3^b$, contrary to prime factorization. Irrational-rotation density gives positive values $\delta(r)$ arbitrarily near zero; hence there are infinitely many strict-record rungs. If $r$ is a rung and $0<t<r$, then $\delta(r)<\delta(t)$, so
\[
\lfloor t\theta+\delta(r)\rfloor=\lfloor t\theta\rfloor.
\]

For $x\in[0,1)$, write
\[
L_t(x)=\lfloor x+t\theta\rfloor,\quad
x_t=\{x+t\theta\},\quad g_t=L_{t+1}-L_t\in\{1,2\}.
\]

The floor-endpoint model has $K_0=0$, $K_n=L_n$, positive integer increments and $K_t\le L_t$. It is in bijection with integer profiles $(j_0,\ldots,j_n)$ satisfying
\[
j_0=j_n=0,\quad j_t\ge0,\quad
j_{t+1}\le j_t+g_t-1,
\]

by $j_t=L_t-K_t$; the inverse increments are $k_t=j_t+g_t-j_{t+1}\ge1$. This is an equality of finite sets, not an asymptotic approximation.

At a rung, take $x=0$. Passing from the source endpoint $\lceil n\theta\rceil$ to $\lfloor n\theta\rfloor$ decreases only the last increment by one. The source cap at $n-1$ implies that this last increment was at least $\lceil n\theta\rceil-\lfloor(n-1)\theta\rfloor\ge2$, so positivity is preserved. Conversely raise the final increment by one. The reward does not use $K_n$, so the transformation preserves it exactly.

For general phase define
\[
w_t=2^{-x_t},\quad B_n(x)=\sum_{t<n}w_t,\quad
E(j)=\sum_{t<n}w_t(1-2^{-j_t}),\quad
Z_n(x,\lambda)=\sum_j e^{-\lambda E(j)}.
\]

Thus the reward is $B_n-E$. Each profile has cost in $[0,n]$, the zero profile is legal, and the profile set is finite. Indeed its inverse positive increments sum to $L_n\le2n$, so it is a subset of finitely many positive compositions. Uniform equidistribution gives
\[
B_n(x)/n\longrightarrow b=\int_0^1 2^{-u}\,du=1/(2\log2)
\]

uniformly in $x$. For completeness, a nonconstant Fourier mode has an irrational-rotation average bounded uniformly in its starting phase by a constant divided by $n$. Trigonometric approximation proves uniform convergence for continuous functions; continuous upper and lower approximants extend it to the bounded functions with finitely many discontinuities used here. This is the uniform phase version of the equidistribution method~\cite{weyl}.

\section{Positive runs and a bounded cutting cost}

A positive run is a maximal consecutive interval $[a,a+\ell-1]$ on which $j>0$. At its first site $j_a=1$ and $g_{a-1}=2$, so $x_a$ is in $[0,\beta)$. Its first cost is at least $2^{-\beta}/2=1/3$. Every subsequent positive site costs at least 1/4. Consequently a run of length $\ell$ has cost at least
\[
1/3+(\ell-1)/4\ge\ell/4.
\]

There are at most $4^{\ell+1}$ shapes for a run of length $\ell$ at a prescribed start. To see this, include its two zero endpoints. Each of the $\ell+1$ depth increments $d$ is at most one and their sum is zero. The numbers $1-d$ are nonnegative integers summing to $\ell+1$. Their number is $\binom{2\ell+1}{\ell}\le4^{\ell+1}$; phase and positivity conditions can only remove shapes.

Setting a cut depth $j_l$ to zero factors the partition exactly. If $j_l>0$, erase the unique positive run containing $l$. Erasure is legal, and the erased profile, interval and shape determine the original profile. There are at most $\ell$ intervals of length $\ell$ containing $l$. Hence, for $\lambda\ge6$,
\[
1\le\frac{Z_{l+m}(x,\lambda)}
 {Z_l(x,\lambda)Z_m(x_l,\lambda)}
\le e^{C_\lambda},\qquad
C_\lambda=\log\!\left(1+\frac{4q}{(1-q)^2}\right),\quad
q=4e^{-\lambda/4}<1.
\]

The cutoff is valid because $4e^{-3/2}<1$. In particular $C_\lambda\le C_6$.

\section{Pressure existence, including phase zero}

For a fixed positive block length $L$, put
\[
p_L(\lambda)=L^{-1}\int_0^1\log Z_L(x,\lambda)\,dx.
\]

The integrand is bounded and piecewise continuous: only finitely many floor breakpoints affect the finite profile set, and within each cell its weights are continuous. Divide $n$ into length-$L$ blocks and a remainder. Iterating the cut inequalities and using uniform equidistribution for rotation by $L\theta$ gives
\[
p_L\le\liminf_n\inf_x n^{-1}\log Z_n(x,\lambda)
\le\limsup_n\sup_x n^{-1}\log Z_n(x,\lambda)
\le p_L+C_6/L.
\]

The remainder has a bounded contribution at fixed $L$: $Z\ge1$ and the number of profiles is bounded in terms of $L$. Since $L$ is arbitrary, the common limit $p(\lambda)$ exists uniformly in x, with
\[
0\le p(\lambda)-p_L(\lambda)\le C_6/L\qquad(\lambda\ge6).
\]

The constant does not depend on $\lambda$. Finite log partition functions are convex, decreasing and 1-Lipschitz after division by length; these properties pass to $p$. The corresponding \textbf{reward} pressure is
\[
P(\lambda)=\lambda b+p(\lambda)\qquad(\lambda\ge6).
\]

Here $Z_n$ is the \emph{defect} partition sum. The source's reward sum is
$Z_n^{\mathrm{rew}}=e^{\lambda B_n(x)}Z_n(x,\lambda)$; the displayed relation
$P(\lambda)=\lambda b+p(\lambda)$ follows from this identity.

This proves the candidate large-tilt convergence assertion only. It does not by itself prove differentiability or convergence for $0\le\lambda<6$.

\section{Matching microcanonical growth by independent finite blocks}

Define $f(\eta)=\inf_{\lambda\ge6}[p(\lambda)+\lambda\eta]$. We claim that, for $0<\eta<10^{-4}$ and $\eta_n\to\eta$, at any sequence of initial phases,
\[
\lim_{n\to\infty}\frac1n\log\#\{j:E(j)\le n\eta_n\}=f(\eta).
\]

The upper bound follows from $\#\{E\le n\eta_n\}\le e^{\lambda n\eta_n}Z_n(x,\lambda)$ for every fixed $\lambda\ge6$ and uniform pressure convergence.

For the lower bound, first ensure an interior finite-block minimizing tilt. Allow singleton runs only at even eligible sites. They are nonadjacent, their density is $\beta$/2, and each costs at most 1/2. Thus
\[
p(6)\ge(\beta/2)\log(1+e^{-3})>1/100.
\]

For example $\beta>1/2$, $e^3<21$ and $\log(1+u)\ge u/(1+u)$ make the lower bound larger than 1/88. Ignoring compatibility of all runs gives
\[
p(40)\le\frac{16e^{-10}}{1-4e^{-10}}<1/1000.
\]

The latter elementary comparison follows from $e^{10}>20000$. These exponential inequalities can also be obtained directly from their Taylor series and tails. Consequently $p(40)+40\eta<p(6)+6\eta$ throughout $0<\eta\le10^{-4}$.

Fix $0<\eta'<\eta$. For all sufficiently large L the same strict comparison holds for $p_L$ by the pressure bracket. Since $p_L\ge0$, the function $p_L(\lambda)+\lambda\eta'$ is coercive and has a minimizer $\lambda_L>6$. Differentiating the finite integral is justified by bounded block cost, and gives
\[
-p_L'(\lambda_L)=\eta'.
\]

On successive length-$L$ blocks, independently sample profiles with probability proportional to $e^{-\lambda_L E}$ at each block's actual starting phase. Concatenate at depth zero and use the zero profile on the remainder. For fixed $L$, equidistribution applied to the bounded piecewise-continuous block mean shows that mean total cost divided by $n$ tends uniformly to $\eta'$. Independence and $E_{\mathrm{block}}\le L$ give total variance $O(nL)$. Chebyshev therefore puts cost divided by $n$ in $[\eta'-d,\eta'+d]$ with probability tending to one, for any $0<d<\min(\eta',\eta-\eta')$. This event lies below $n\eta_n$ for all large $n$.

Undoing the Gibbs weights bounds its number of profiles below by
\[
\Pr(\text{event})\,
 \prod_{\text{blocks}}Z_L(x_{\text{block}},\lambda_L)
 \exp[\lambda_L n(\eta'-d)].
\]

Its logarithm/n has lower limit at least $p_L(\lambda_L)+\lambda_L(\eta'-d)$. First let d decrease to zero. Then let L grow: the uniform bracket $|p-p_L|\le C_6/L$ implies convergence of the minimizing values to $f(\eta')$. Finally let $\eta'$ increase to $\eta$. The function f is finite, nonnegative and concave, so it is continuous at every $\eta>0$ in this open interval; compare the general continuity theory in~\cite{rockafellar}. This proves the claimed lower bound. No derivative of the infinite pressure was assumed.

At source rungs, $E\le\epsilon B_r$ corresponds to $\eta_r=\epsilon B_r/r\to\epsilon b$. Therefore
\[
\rho(1-\epsilon)=f(\epsilon b)/\log10
\]

in the indicated near-one interval. At $\epsilon=0$ every positive depth has strictly positive cost, so the zero profile is the unique qualifying word.

\section{An explicit large-tilt sandwich}

At an interior site $a$ with $1\le a<n$, a singleton is legal precisely when $x_a$ is in $[0,\beta)$; its activity is $v_a=e^{-\lambda 2^{-x_a}/2}$, and is zero otherwise. Its mean activity is
\[
I(\lambda)=\int_0^\beta e^{-\lambda 2^{-u}/2}\,du
 =\frac1{\log2}\int_{1/3}^{1/2}\frac{e^{-\lambda y}}y\,dy.
\]

For an upper bound, ignore all incompatibility between positive runs and use $\log(1+v)\le v$. Singleton contributions have mean I; runs of length at least two have total activity per possible start bounded by
\[
\sum_{\ell\ge2}4^{\ell+1}
 e^{-\lambda(1/3+(\ell-1)/4)}
 =\frac{64e^{-7\lambda/12}}{1-4e^{-\lambda/4}}.
\]

For a lower bound keep only nonadjacent singleton runs. Their finite partition recurrence is $H_m=H_{m-1}+v_mH_{m-2}$, with empty-boundary values equal to one. The ratio $R_m=H_m/H_{m-1}$ lies between $1+v_m/(1+v_{m-1})$ and $1+v_m$. Since $\log(1+u)\ge u-u^2/2$ for $u\ge0$,
\[
\log R_m\ge v_m-v_mv_{m-1}-v_m^2/2.
\]

Every nonzero $v_m$ is at most $e^{-\lambda/3}$. Summing and phase averaging yields
\[
I(\lambda)-\tfrac32e^{-2\lambda/3}
\le p(\lambda)
\le I(\lambda)+\frac{64e^{-7\lambda/12}}{1-4e^{-\lambda/4}}.
\]

Substitute $y=1/3+z$ in $I$. The factor $1/(1/3+z)=3+O(z)$ on $[0,1/6]$, and its bounded derivative controls the error after integration against $e^{-\lambda z}$. The missing tail beyond $1/6$ is exponentially small. Thus
\[
I(\lambda)=\frac{3e^{-\lambda/3}}{\lambda\log2}
 (1+O(1/\lambda)),\qquad
p(\lambda)\sim\frac{3e^{-\lambda/3}}{\lambda\log2}.
\]

These are estimates in the untruncated profile model, not a finite-matrix extrapolation. Also $p(\lambda)>0$ at every finite $\lambda\ge6$ by the even-singleton construction, and $p(\lambda)\to0$ by the upper bound.

\section{The asymptotic Legendre minimum}

The minimizing tilts for $f(\eta)$ escape to infinity as $\eta$ decreases to zero. Indeed $f(\eta)\to0$ by first choosing a large fixed tilt and then decreasing $\eta$; whereas a bounded tilt interval has a positive lower bound for p. For every fixed $0<d<1$, the preceding asymptotic sandwiches p, for sufficiently large tilt, between
\[
(1\pm d)\frac{3e^{-\lambda/3}}{\lambda\log2}.
\]

The minimizers of these two comparison objectives also escape to infinity, so the sandwich applies to their infima for small $\eta$. Put $t=\lambda/3$ and $C=(1\pm d)/\log2$. The comparison objective is $Ce^{-t}/t+3\eta t$. Its unique positive minimizer satisfies
\[
3\eta=C e^{-t}(1+1/t)/t.
\]

Writing $A=\log(1/\eta)$ and taking logarithms gives
\[
t+\log t=A+\log(C/3)+\log(1+1/t),\qquad
 t=A-\log A+\log(C/3)+o(1).
\]

At the minimizer the first term is $3\eta/(1+1/t)$. Therefore the comparison minimum equals
\[
3\eta[A-\log A+1+\log(C/3)+o(1)].
\]

Taking $\eta$ to zero first and then d to zero squeezes the actual minimum to
\[
f(\eta)=3\eta\left[\log(1/\eta)-\log\log(1/\eta)
 +1-\log(3\log2)+o(1)\right].
\]

Substitute $\eta=\epsilon/(2\log2)$ and divide by $\log10$. The constants combine as $\log(2\log2)-\log(3\log2)=\log(2/3)$, giving the expansion in Section 1.

\section{Computation, formal coverage and remaining work}

The \href{../../certificates/external_problems/bridge.py}{exact finite bridge oracle} and its recorded data check literal finite words against depth profiles, with rational caps and rewards. They include nonrungs; a finite normalization witness at n=2 has words (1,3),(2,2), rewards 5/3,7/3, and phase-zero $B_2=5/3$. Both meet that fixed threshold, while only the latter meets the actual maximum. This witnesses the difference of definitions; it does not by itself refute an asymptotic limit. The scope of each executed check is recorded separately.

\begin{center}\small\begin{tabular}{@{}>{\raggedright\arraybackslash}p{0.17\textwidth}>{\raggedright\arraybackslash}p{0.34\textwidth}>{\raggedright\arraybackslash}p{0.41\textwidth}@{}}
\toprule

Target & This draft supplies & Still required for the release goal\\[0.6em]

\midrule

Near-one rung M4 & Written entropy-existence and variational argument & Frozen full formal statement, end-to-end proof, certified entropy table, complete assurance\\[0.6em]

R9 & Written leading and second-order expansion & Full formal analytic bounds and limits, statement matching, fresh and independent checker replay\\[0.6em]

R8 & Written convergence for $\lambda\ge6$ & All-tilt convergence and differentiability, including critical tilt, independently scrutinized and fully formalized\\[0.6em]

Fixed-reference all-length endpoint & Appendix~\ref{app:all-length}: explicit nonconvergence construction at c=1, with 23 formal source proofs & Full independent assurance for that separate package; no intended rung M4 or phase-adapted conclusion\\[0.6em]

\bottomrule\end{tabular}\end{center}

Broader all-phase, all-tilt and critical Gibbs-law arguments are research candidates. They are not premises of the near-extremal derivation above and are not accepted theorems in this release. A separate critical argument would need a rigorous stationary finite-depth law, preservation of its conditional Gibbs identities, same-law excursion bounds and return-gap contradiction; none is supplied by the finite oracle or supporting Lean library.

The outstanding trust gates include independent review of the full formal statements against the pinned source, full analytic formalization and transitive axiom audit, fresh proof replay and independently implemented checker coverage of these actual headline claims. Supporting-theorem checks must remain labeled as such. Closed review gates for supporting arithmetic do not close these gaps.

\appendix
\section{The elementary analytic tools and explicit bounds}
For an irrational rotation by $\omega$ and any nonzero integer $h$, the
geometric-sum identity gives
\[
\left|\frac1n\sum_{t<n}e^{2\pi ih(x+t\omega)}\right|
 \le \frac{2}{n|1-e^{2\pi ih\omega}|},
\]
uniformly in the initial phase $x$. Trigonometric-polynomial approximation
therefore proves uniform averaging for continuous periodic functions.
A bounded Riemann-integrable function can be bracketed above and below by
continuous periodic functions whose integrals differ arbitrarily little:
use Darboux step bounds and continuous ramps in small neighborhoods of the
finitely many partition endpoints. This proves the uniform version used
for the finite-block integrands, their mean costs and $u\mapsto2^{-u}$.
The argument applies to $L\theta$ as well, since every positive integer
multiple of an irrational number is irrational.

For independent bounded block costs, variance adds. With at most $n/L+1$
blocks of cost in $[0,L]$, the total variance is at most $(n/L+1)L^2$.
Markov's inequality applied to the squared centered cost gives the stated
Chebyshev estimate. Independence is chosen in this finite counting proof;
it is not an assertion about the exponents of an integer Collatz trajectory.

The numerical comparisons in Section~5 can be checked by rational Taylor
bounds, with no floating-point premise. In particular,
\[
 e^3\le\sum_{j=0}^{8}\frac{3^j}{j!}
       +\frac{3^9/9!}{1-3/10}<21,
 \qquad
 e^{10}>\sum_{j=0}^{14}\frac{10^j}{j!}>20000.
\]
For the upper tail of $e^3$, consecutive omitted terms have ratio at most
$3/10$. Also $3^2>2^3$ proves $\beta>1/2$, and integration of
$(1+t)^{-1}\ge(1+u)^{-1}$ on $0\le t\le u$ proves
$\log(1+u)\ge u/(1+u)$. These give the explicit strict pressure comparison
used to keep the finite-block minimizing tilt inside $(6,\infty)$.

The continuity fact used in the microcanonical lower bound is local: a finite
concave function on an open interval is bounded between its neighboring secant
slopes on every compact subinterval, and is therefore continuous there. Thus
no differentiability of the infinite-volume pressure is used.

\section{The fixed-reference all-length extension at c=1}\label{app:all-length}
This appendix concerns a separate endpoint distinction. Define $N_{\rm all}(1,r)$
using the positive ordered source words at every positive length $r$, closing sum
$s=\lceil r\theta\rceil$, shifted interior caps
$K_t\le\lfloor t\theta+\delta(r)\rfloor$, and the fixed phase-zero threshold
$\Phi\ge B_r$. Do not replace that threshold by a phase-adapted maximum.
The isolated M4 wording admits this all-length reading, but its status paragraph
and detailed continuation point to rungs. The following is not a refutation of
the intended rung-only question, the source's phase-adapted $N^*$, or a theorem
at $c<1$. No extra integer-cycle or rotation-quotient condition is imposed.

\begin{theorem}[All-length endpoint distinction]\label{thm:all-length}
For the actual slope $\theta=\log3/\log2$, put
$a_r=r^{-1}\log_{10}N_{\rm all}(1,r)$ for $r\ge1$ and $a_0=0$. Then
\[
 \liminf_{r\to\infty}a_r=0,\qquad
 \limsup_{r\to\infty}a_r\ge\frac{\log_{10}2}{8}>0.
\]
In particular the all-length limit does not exist. More explicitly,
\[
 N_{\rm all}(1,r)\ge2^{\lfloor r/8\rfloor}
 \quad\hbox{if }r\ge8\hbox{ and }\delta(r)\ge\tfrac12,
\]
whereas $N_{\rm all}(1,r)=1$ at every positive strict rung.
\end{theorem}
\begin{proof}
The exact comparisons $2^3<3^2$ and $3^3<2^5$ give
$3/2<\theta<5/3$, hence $1/2<\beta=\theta-1<2/3$.
For every phase $y\in[0,1)$ the three phases
$y,\{y+\beta\},\{y+2\beta\}$ include one in each half-circle.
For the high half, if $y<1/2$ and its successor wraps, that successor is
$z<\beta-1/2<1/6$, so $z+\beta\in(1/2,5/6)$ without wrapping; if it does
not wrap it is already above $1/2$. For the low half, start at $y\ge1/2$.
Its successor is $z<\beta$; if $z\ge1/2$, the next phase wraps and is
$z+\beta-1<2\beta-1<1/3$. Thus among indices $0\le t<r$ at least
$\lfloor r/3\rfloor$ phases $\{t\theta\}$ are high. The low-half fact in
arbitrarily late triples, together with irrationality excluding phase zero
at positive indices, also gives unbounded $r$ with $\delta(r)\ge1/2$.

Use the canonical prefixes $C_t=\lfloor t\theta+\delta(r)\rfloor$ for
$0\le t\le r$. They start at zero, end at $s$, and their successive letters
are 1 or 2. Two consecutive letters cannot both be 1, since their sum is at
least $\lfloor2\theta\rfloor\ge3$. The number of 1 letters is
\[
 2r-s\ge2r-\lceil5r/3\rceil=\lfloor r/3\rfloor.
\]
Except possibly for the first letter, each 1 is preceded by 2. These 21 pairs
are disjoint. For $r\ge8$, write $m=\lfloor r/8\rfloor\ge1$; then
$\lfloor r/3\rfloor\ge2m$ and there are at least $m$ such pairs.

Select $m$ pairs. Toggling any subset 21 to 12 preserves positivity and the
closing sum, and lowers only each pair's intervening interior prefix by one.
All source caps remain satisfied. Distinct subsets give distinct prefixes,
hence $2^m$ distinct ordered words.
When $\delta(r)\ge1/2$, every high phase has
$C_t=\lfloor t\theta\rfloor+1$, so the canonical excess obeys
\[
 \Phi(C)-B_r\ge
 \sum_{\substack{t<r\\\{t\theta\}\ge1/2}}2^{-\{t\theta\}}
 \ge\tfrac12\lfloor r/3\rfloor\ge m.
\]
A toggle at the prefix $t$ loses exactly
$2^{C_t-1}/3^t=2^{C_t-t\theta-1}<1$, because
$C_t-t\theta\le\delta(r)<1$. Disjoint toggles affect different reward
terms. Every selected subset therefore loses less than $m$, proving the
exponential lower bound at the same fixed reference. The identity
$3^t=2^{t\theta}$ uses the actual slope; interval bounds on a generic slope
alone would not justify this reward calculation.

For completeness the counts are finite and positive at every positive length.
The canonical word is admissible and has reward at least $B_r$.
A selected word is an ordered positive composition of $s$; the complete list
of its letters determines it. There are $2^{s-1}$ such compositions, and
$s\le2r$, giving
$1\le N_{\rm all}(1,r)\le2^{2r}$ and
$0\le a_r\le2\log_{10}2$. At length zero we use the formal convention of an empty source population, and
$a_0=0$ is only a harmless sequence convention. Thus the real
liminf and limsup are finite and well posed.

At a strict rung the shifted caps equal $\lfloor t\theta\rfloor$ for every
interior index. Every reward term is bounded by its phase-zero reference
term, and equality of the sum forces equality of every interior prefix.
The closing sum determines the last letter, giving exactly one qualifying
word. Actual positive strict rungs are unbounded, so nonnegativity and this
subsequence give $\liminf a_r=0$. On the unbounded high-defect lengths,
\[
 a_r\ge\frac{\lfloor r/8\rfloor}{r}\log_{10}2,
 \qquad \frac{\lfloor r/8\rfloor}{r}\longrightarrow\frac18.
\]
This proves the positive lower bound on the limsup and excludes a real limit.
\end{proof}

\paragraph{Separate formal scope.}
The public \path{lean/external_unrestricted/} package implements 23 frozen
statements for this exact distinction. A fresh ordinary build matched all 23
proposition bodies and 13 model definitions and audited 504 theorem/helper
closures using only the standard three axioms. Its complete source-bound
replay, trusted comparator, six checker modes and hostile-control gates are
\textbf{pending at this manuscript snapshot}. The
\href{../../lean/external_unrestricted/ARGUMENT.md}{full companion}
includes every finite-population, swap and limit argument and its formal map.
The underlying finite bridge package has separate successful source-pinned
receipts. These formalizations use mathlib~\cite{mathlib}, at the exact
revision fixed in each public lockfile. They do not formalize the main
rung-pressure, microcanonical or R9 claims of this research draft, and do
not supply novelty certification or independent human expert refereeing.

\begin{thebibliography}{9}
\bibitem{polson}
N.~Polson and V.~Sokolov.
\emph{On the Collatz Conjecture: Martingales, Momentum and Continued Fractions}.
Edition~1, version~1.4.0, 2026.
\href{https://doi.org/10.5281/zenodo.24113412}{doi:10.5281/zenodo.24113412}.
The labels \texttt{rrt:sec:m4}, \texttt{rrt:prob:r8} and
\texttt{rrt:prob:r9} refer to the
\href{https://github.com/VadimSokolov/collatz/blob/5840497a23fc0d64127b850aea68862d7c93f2b7/book/chapters/ch10h_ratchets.tex}{pinned ratchets chapter}.
The unrestricted M4 wording appears in the
\href{https://github.com/VadimSokolov/collatz/blob/5840497a23fc0d64127b850aea68862d7c93f2b7/book/chapters/ch10f_missed.tex}{pinned missed-problems chapter}.
\bibitem{mathlib}
The mathlib Community. The Lean mathematical library.
\emph{Proceedings of CPP 2020} (2020), 367--381.
\href{https://doi.org/10.1145/3372885.3373824}{doi:10.1145/3372885.3373824};
\href{https://arxiv.org/abs/1910.09336}{author manuscript}.

\bibitem{halbeisen-hungerbuhler}
L.~Halbeisen and N.~Hungerb\"uhler.
Optimal bounds for the length of rational Collatz cycles.
\emph{Acta Arithmetica} \textbf{78}(3) (1997), 227--239.
\href{https://doi.org/10.4064/aa-78-3-227-239}{doi:10.4064/aa-78-3-227-239}.

\bibitem{weyl}
H.~Weyl. \"Uber die Gleichverteilung von Zahlen mod.~Eins.
\emph{Mathematische Annalen} \textbf{77} (1916), 313--352.
\href{https://doi.org/10.1007/BF01475864}{doi:10.1007/BF01475864}.

\bibitem{giacomin}
G.~Giacomin. Renewal sequences, disordered potentials, and pinning phenomena.
Lecture notes (2008), Sections~1.2 and~3.1.
\href{https://arxiv.org/abs/0807.4285}{arXiv:0807.4285}.

\bibitem{rockafellar}
R.~T.~Rockafellar. \emph{Convex Analysis}.
Princeton University Press, 1970, Parts~II--III.
\href{https://doi.org/10.1515/9781400873173}{doi:10.1515/9781400873173}.

\bibitem{dlmf}
NIST. \emph{Digital Library of Mathematical Functions}, Section~2.3(ii),
Watson's lemma, equation~(2.3.8).
\href{https://dlmf.nist.gov/2.3#ii}{dlmf.nist.gov/2.3(ii)}.
Accessed 27 November 2026.
\end{thebibliography}
\end{document}
